\begin{ex}
	Rút gọn biểu thức $P=\frac{\sqrt{x}}{\sqrt{x}+1}+\frac{1}{\sqrt{x}-1}+\frac{2\sqrt{x}}{1-x}$ (với $x \ge 0, x \ne 1$)
	\loigiai{
		Điều kiện: $x \ge 0, x \ne 1$.
		Ta có:
		\begin{align*}
			P &= \frac{\sqrt{x}}{\sqrt{x}+1} + \frac{1}{\sqrt{x}-1} + \frac{2\sqrt{x}}{1-x} \\
			&= \frac{\sqrt{x}(\sqrt{x}-1)}{(\sqrt{x}+1)(\sqrt{x}-1)} + \frac{1(\sqrt{x}+1)}{(\sqrt{x}-1)(\sqrt{x}+1)} - \frac{2\sqrt{x}}{x-1} \\
			&= \frac{x-\sqrt{x}}{x-1} + \frac{\sqrt{x}+1}{x-1} - \frac{2\sqrt{x}}{x-1} \\
			&= \frac{x-\sqrt{x}+\sqrt{x}+1-2\sqrt{x}}{x-1} \\
			&= \frac{x-2\sqrt{x}+1}{x-1} \\
			&= \frac{(\sqrt{x}-1)^2}{(\sqrt{x}-1)(\sqrt{x}+1)} \\
			&= \frac{\sqrt{x}-1}{\sqrt{x}+1}
		\end{align*}
		Vậy, với $x \ge 0, x \ne 1$ thì $P = \frac{\sqrt{x}-1}{\sqrt{x}+1}$.
	}
\end{ex}

\begin{ex} 
	Gọi $x_1, x_2$ là hai nghiệm của phương trình: $x^2 - 4x - 7 = 0$. Không giải phương trình, tính giá trị của biểu thức $T = \frac{x_1}{x_2} + \frac{x_2}{x_1} - 2$
	\loigiai{
		Phương trình đã cho là $x^2 - 4x - 7 = 0$.
		Ta biến đổi biểu thức $T$:
		\begin{align*}
			T &= \frac{x_1}{x_2} + \frac{x_2}{x_1} - 2 \\
			&= \frac{x_1^2 + x_2^2}{x_1 x_2} - 2 \\
			&= \frac{(x_1 + x_2)^2 - 2x_1 x_2}{x_1 x_2} - 2
		\end{align*}
		Thay giá trị từ hệ thức Vi-ét vào biểu thức $T$:
		\begin{align*}
			T &= \frac{(4)^2 - 2(-7)}{-7} - 2 \\
			&= \frac{16 + 14}{-7} - 2 \\
			&= \frac{30}{-7} - 2 \\
			&= -\frac{30}{7} - \frac{14}{7} \\
			&= -\frac{44}{7}
		\end{align*}
		Vậy, giá trị của biểu thức $T = -\frac{44}{7}$.
	}
\end{ex}

\begin{ex} % Câu 11
	Giải hệ phương trình:
	$ \begin{cases} 2x-y=3 \\ x-2y=3 \end{cases} $
	\loigiai{
		Ta có hệ phương trình:
		$ \begin{cases} 2x-y=3 \quad &(1) \\ x-2y=3 \quad &(2) \end{cases} $
		Từ phương trình (1), ta có $y = 2x - 3$.
		Thế $y = 2x - 3$ vào phương trình (2), ta được:
		\begin{align*}
			x - 2(2x-3) &= 3 \\
			x - 4x + 6 &= 3 \\
			-3x &= 3 - 6 \\
			-3x &= -3 \\
			x &= 1
		\end{align*}
		Thay $x=1$ vào biểu thức $y = 2x - 3$:
		\begin{align*}
			y &= 2(1) - 3 \\
			y &= 2 - 3 \\
			y &= -1
		\end{align*}
		Vậy, nghiệm của hệ phương trình là $(x; y) = (1; -1)$.
	}
\end{ex}
\begin{ex} % Câu a - Xác suất
	Một hộp có 12 thẻ cùng loại, mỗi thẻ được ghi một trong các số $1; 2; 3; \dots; 12$; hai thẻ khác nhau thì ghi hai số khác nhau. Xét phép thử “Thẻ được rút ra ghi số chia hết cho 3”. Tính xác suất của biến cố M?
	\loigiai{
		Không gian mẫu $\Omega = \{1; 2; 3; 4; 5; 6; 7; 8; 9; 10; 11; 12\}$.\\
		Số phần tử của không gian mẫu là $n(\Omega) = 12$.\\
		Biến cố M: “Thẻ được rút ra ghi số chia hết cho 3”.\\
		Các phần tử của biến cố M là $M = \{3; 6; 9; 12\}$.\\
		Số phần tử của biến cố M là $n(M) = 4$.\\
		Xác suất của biến cố M là:\\
		$P(M) = \frac{n(M)}{n(\Omega)} = \frac{4}{12} = \frac{1}{3}$.
	}
\end{ex}
\begin{ex} % Bài toán mua sắm
	Để chuẩn bị cho năm học mới, bạn Đỉnh đến một cửa hàng để mua vở và bút. Hôm đó, cửa hàng mừng khai trương nên giảm giá 10\% cho vở và 15\% cho bút. Vì vậy, tổng số tiền vở và bút mà Đỉnh đã trả là 176 nghìn đồng. Chị bán hàng cho biết, nếu cửa hàng không giảm giá thì tổng số tiền mà Đỉnh phải trả sẽ là 200 nghìn đồng. Em hãy tìm số tiền mua vở, số tiền mua bút mà Đỉnh phải trả nếu cửa hàng không giảm giá.
	\loigiai{
		Gọi $x$ (nghìn đồng) là số tiền mua vở ban đầu (khi chưa giảm giá).
		Gọi $y$ (nghìn đồng) là số tiền mua bút ban đầu (khi chưa giảm giá).
		Điều kiện: $x > 0, y > 0$.
		
		Theo đề bài, nếu cửa hàng không giảm giá thì tổng số tiền mà Đỉnh phải trả sẽ là 200 nghìn đồng, nên ta có phương trình:
		$x + y = 200 \quad (1)$
		
		Cửa hàng giảm giá 10\% cho vở, nên số tiền mua vở sau khi giảm giá là:
		$x - 0,1x = 0,9x$ (nghìn đồng).
		Cửa hàng giảm giá 15\% cho bút, nên số tiền mua bút sau khi giảm giá là:
		$y - 0,15y = 0,85y$ (nghìn đồng).
		
		Tổng số tiền vở và bút mà Đỉnh đã trả sau khi giảm giá là 176 nghìn đồng, nên ta có phương trình:
		$0,9x + 0,85y = 176 \quad (2)$
		
		Từ (1) và (2), ta có hệ phương trình:
		$ \begin{cases} x + y = 200 \\ 0,9x + 0,85y = 176 \end{cases} $
		
		Giải hệ phương trình:
		Từ phương trình (1), suy ra $y = 200 - x$.
		Thế $y = 200 - x$ vào phương trình (2):
		\begin{align*}
			0,9x + 0,85(200 - x) &= 176 \\
			0,9x + 170 - 0,85x &= 176 \\
			0,05x &= 176 - 170 \\
			0,05x &= 6 \\
			x &= \frac{6}{0,05} \\
			x &= 120
		\end{align*}
		Thay $x = 120$ vào $y = 200 - x$:
		\begin{align*}
			y &= 200 - 120 \\
			y &= 80
		\end{align*}
		Vậy, nếu cửa hàng không giảm giá, số tiền mua vở là 120 nghìn đồng và số tiền mua bút là 80 nghìn đồng.
	}
\end{ex}
\begin{ex} % Câu 13
	Cho tam giác ABC nhọn (AB < AC) nội tiếp đường tròn (O;R). Các đường cao AK, BE, CD cắt nhau tại H.
	\begin{enumerate}[a)]
		\item Chứng minh bốn điểm B, D, E, C cùng thuộc một đường tròn.
		\item Gọi N là giao điểm của DC và KE, P là hình chiếu vuông góc của C lên đường kính AM. Khi A, B cố định, C thay đổi nhưng vẫn thỏa mãn bài toán, chứng minh NH.CD = DH.NC và xác định vị trí của C để chu vi tam giác OCP lớn nhất.
	\end{enumerate}
\end{ex}
